\scanpage{33}
% archive foliation: 32
$\Phi(F_1,-F_n)$, où le $F_i$ est un faisceau coh. sur un $=X\otimes_kK$, \inserted{$K$ alg. clos,} $\operatorname{cl}_kF_i\in A_i$.

D'ailleurs, si on \uline{part} d'\deleted{\unsure{autres}} faisceaux \deleted{$K_i$} $F_i$ sur des extensions \uline{distinctes} $K_i/k$, $F_i$ de classe $\xi_i\in\unsure{\Phi|\mathrm{Cl}}(X/k)$, on obtient $\Phi(\xi_1,-\xi_n)$ en prenant de toutes les façons possibles une extension composée des $K_i$ sur $k$, et sa clôture algébrique \deleted{$K_\alpha$} $L_\alpha$, et l'on \unsure{réunit|somme} les ens de classes dans $\unsure{\Phi|\mathrm{Cl}}(X/k)$ qui correspondent à
\[
\deleted{\Phi(F_1^{(\alpha)},F_n^{(\alpha)})}\quad
\Phi(F_1\otimes_kL_\alpha,-,F_n\otimes_kL_\alpha).
\]
\margin{par le fait que les $F_i$ sont \unsure{sur|les} les \unsure{fibres} des $\Rightarrow$ \unsure{schémas|faisceaux}, $\checkmark$ \illegible{2} \unsure{qu'on a des} $\unsure{X_i-\alpha}$ \unsure{réunion}}

Si $\Phi$ est défini aussi pour des faisceaux sur un $X_K$, sans supposer $K$ alg. clos, avec validité correspondante pour (i) et (ii), alors on n'a pas besoin dans les applications précédentes de \deleted{\unsure{prendre}}\inserted{supposer les $K_i$ alg. clos}, et de passer à la clôture algébrique des extensions composées.
